\chapter{Architecture du lab UrbanUpC}
\label{ch:lab-archi}

Ce chapitre décrit comment le lab est monté sur Azure. L'objectif
en construisant ce lab était d'avoir un SI suffisamment réaliste
pour avoir quelque chose à homologuer~: un domaine Active
Directory, des applications web exposées, un SIEM qui détecte les
attaques. Pas un jouet, mais pas non plus une infra de production.

\section{Plateforme et budget}

Tout tourne sur \textbf{Microsoft Azure}, région
\texttt{swedencentral}. Le compte utilisé est un \emph{Azure for
Students} (100\,\euro\ de crédit), ce qui force à être économe~:
les VMs sont en taille \texttt{Standard\_B2s\_v2} (2 vCPU,
8\,Go~RAM), et un auto-shutdown à 22h~UTC arrête tout pour la
nuit.

Le \emph{Resource Group} unique \texttt{RG-PFE-SOC} contient
toute l'infra~: trois VMs, un \emph{Virtual Network} segmenté en
trois subnets, trois \emph{Network Security Groups}, et les
disques associés.

\section{Vue d'ensemble}

\begin{figure}[h]
\centering
\begin{small}
\begin{verbatim}
                    Internet
                       |
                 [IPs publiques]
                       |
+=========[vnet-pfe 10.0.0.0/16]==========+
|                                          |
|  snet-dmz (10.0.1.0/24)     nsg-dmz     |
|  +-------------------+                   |
|  | vm-web01          | Ubuntu 22.04 LTS |
|  | 10.0.1.10         | Docker + apps    |
|  | :80 :443 :8888    | + agent Wazuh    |
|  +---------+---------+                   |
|            X (Deny-DMZ-to-LAN)           |
|                                          |
|  snet-lan (10.0.2.0/24)     nsg-lan     |
|  +-------------------+                   |
|  | vm-dc01           | Windows Srv 2022 |
|  | 10.0.2.10         | AD corpnet.local |
|  | :3389 (RDP)       | + agent Wazuh    |
|  +---------+---------+                   |
|            |                             |
|  snet-mgmt (10.0.3.0/24)    nsg-mgmt    |
|  +-------------------+                   |
|  | vm-siem01         | Ubuntu 22.04 LTS |
|  | 10.0.3.10         | Wazuh all-in-one |
|  | :443 :1514 :1515  | + dashboard      |
|  +-------------------+                   |
+==========================================+
\end{verbatim}
\end{small}
\caption{Schéma réseau du lab UrbanUpC sur Azure.}
\label{fig:archi-lab}
\end{figure}

\section{Les trois VMs}

\begin{table}[h]
\centering
\caption{Les trois machines virtuelles du lab.}
\label{tab:vms}
\small
\begin{tabularx}{\textwidth}{lXXXX}
\toprule
\textbf{VM} & \textbf{OS} & \textbf{Subnet} & \textbf{IP privée} & \textbf{Rôle} \\
\midrule
vm-web01 & Ubuntu 22.04 LTS & snet-dmz & 10.0.1.10 & Serveur web exposé sur Internet, héberge les applications PHP \\
vm-dc01 & Windows Server 2022 Datacenter & snet-lan & 10.0.2.10 & Contrôleur de domaine Active Directory \texttt{corpnet.local} \\
vm-siem01 & Ubuntu 22.04 LTS & snet-mgmt & 10.0.3.10 & SIEM Wazuh all-in-one (manager + indexer + dashboard) \\
\bottomrule
\end{tabularx}
\end{table}

L'authentification est par clé SSH pour les deux Ubuntu, par mot
de passe pour le DC Windows (contrainte du provisionnement Azure).
Toutes les VMs ont un agent Wazuh installé qui remonte les logs
vers \texttt{vm-siem01}.

\section{Segmentation réseau}

Trois subnets, trois NSG. La logique~:

\begin{itemize}
  \item \textbf{snet-dmz} (10.0.1.0/24)~: zone exposée à Internet.
    Le serveur web y vit. Règle de sortie clé~: \texttt{Deny-DMZ-to-LAN}
    interdit toute communication depuis la DMZ vers le LAN. Si
    un attaquant compromet le web, il ne peut pas pivoter vers
    l'AD directement.
  \item \textbf{snet-lan} (10.0.2.0/24)~: réseau interne, où vit
    l'AD. Pas d'accès Internet (\texttt{Deny-LAN-to-Internet}),
    comme dans un vrai SI. Le DC ne peut sortir que vers le SIEM
    pour les logs.
  \item \textbf{snet-mgmt} (10.0.3.0/24)~: management. Seuls les
    agents Wazuh (ports 1514/1515) et l'administrateur (SSH,
    dashboard sur 443/55000) y accèdent.
\end{itemize}

C'est la segmentation classique DMZ / LAN / MGMT, à laquelle on a
ajouté l'isolation \texttt{Deny-DMZ-to-LAN} qui est ce qui rend
le lab intéressant pour tester un scénario de pivot.

\section{Règles NSG principales}

\begin{table}[h]
\centering
\caption{Règles NSG les plus structurantes (extrait).}
\label{tab:nsg}
\small
\begin{tabularx}{\textwidth}{llcllX}
\toprule
\textbf{NSG} & \textbf{Règle} & \textbf{Pri} & \textbf{Sens} & \textbf{Action} & \textbf{Description} \\
\midrule
nsg-dmz & Allow-HTTP & 100 & In & Allow & 80/443 depuis Internet \\
nsg-dmz & Allow-SSH-Admin & 110 & In & Allow & 22 depuis IP admin uniquement \\
nsg-dmz & Deny-DMZ-to-LAN & 200 & Out & Deny & Tout vers 10.0.2.0/24 \\
nsg-lan & Allow-RDP-Admin & 100 & In & Allow & 3389 depuis IP admin \\
nsg-lan & Deny-LAN-to-Internet & 200 & Out & Deny & LAN ne sort pas sur Internet \\
nsg-mgmt & Allow-Wazuh-Agents & 110 & In & Allow & 1514/1515 depuis vnet \\
\bottomrule
\end{tabularx}
\end{table}

L'IP admin est mise à jour à la main quand elle change, via
\texttt{az network nsg rule update}. C'est pénible mais
acceptable pour un lab.

\section{Provisionnement reproductible}

Tout l'environnement peut être recréé à partir de quatre scripts~:

\begin{itemize}
  \item \texttt{infra/azure-deploy.sh}~: crée le Resource Group,
    le VNet, les NSG, les trois VMs. Idempotent.
  \item \texttt{infra/vm-web01-setup.sh}~: installe Docker et la
    stack applicative sur le serveur web.
  \item \texttt{infra/vm-siem01-setup.sh}~: installe Wazuh
    all-in-one (manager + indexer + dashboard) sur le SIEM.
  \item \texttt{infra/vm-dc01-setup.ps1}~: installe AD~DS et
    crée le domaine \texttt{corpnet.local} sur le contrôleur.
\end{itemize}

Cette reproductibilité est importante~: c'est elle qui permet à
l'outil HOMO-CI d'avoir un état attendu à comparer à l'état réel.

\begin{takeaway}
Le lab tient sur 3 VMs Azure très modestes, segmentées en
DMZ/LAN/MGMT. C'est volontairement minimaliste, pour rester
dans le crédit étudiant et garder l'infra compréhensible.
\end{takeaway}
